\documentclass[11pt]{article}
\usepackage[margin=1in]{geometry}
\usepackage{parskip}
\usepackage{hyperref}

\title{Lecture Notes Script (Slide-by-Slide): Lecture 2}
\author{Hui-Jun Chen}
\date{\today}

\begin{document}
\maketitle

\section*{Lecture 2 (RBC $\rightarrow$ Money Market $\rightarrow$ IS--LM $\rightarrow$ AD)}

\begin{enumerate}
\item \textbf{Today: why we move from RBC to IS--LM}
\begin{itemize}
\item Last semester RBC pinned down \textbf{real allocations}; this semester we care about \textbf{inflation and the price level}.
\item To talk about $P$ and inflation, we add a market where nominal variables matter: \textbf{money demand / money market}.
\item Today’s bridge model: \textbf{IS (goods)} + \textbf{LM (money)}.
\end{itemize}

\item \textbf{RBC in one slide: what we learned (and what we did not)}
\begin{itemize}
\item RBC gives a disciplined \textbf{real core}: technology/labor determine $Y$, Euler pins $r$.
\item RBC alone does not explain \textbf{nominal rates, money demand, or the price level}.
\item We keep the real logic and put the nominal side on the board.
\end{itemize}

\item \textbf{Competitive Equilibrium}
\begin{itemize}
\item RBC equilibrium map: labor market $\Rightarrow (N,w)$; production $\Rightarrow Y$; intertemporal market $\Rightarrow r$.
\item Use this as the \textbf{long-run benchmark}.
\item Next: add money market.
\end{itemize}

\item \textbf{A clean map: same economy, more markets on the board}
\begin{itemize}
\item Macro equilibrium = clearing conditions: \textbf{goods market + money market + policy}.
\item IS--LM is a \textbf{representation} highlighting demand + nominal variables.
\item Start with money side.
\end{itemize}

\item \textbf{Step 1: Money demand (why people hold money)}
\begin{itemize}
\item Money provides liquidity; holding it has opportunity cost (interest).
\item $M^d/P = L(Y,i)$: $Y\uparrow \Rightarrow L\uparrow$; $i\uparrow \Rightarrow L\downarrow$.
\item Next: what shifts money demand.
\end{itemize}

\item \textbf{Money demand shifts I: higher price level means higher nominal money demand}
\begin{itemize}
\item Real balances matter: $M/P$. If $P\uparrow$, same real balances require more nominal money.
\item For given $(Y,i)$, nominal money demand rises.
\item Next: income shift.
\end{itemize}

\item \textbf{Money demand shifts II: higher income means more transactions}
\begin{itemize}
\item $Y\uparrow$ increases transactions $\Rightarrow$ money demand rises.
\item This makes LM slope upward later.
\item Next: money supply.
\end{itemize}

\item \textbf{Money supply: the central bank side (first pass)}
\begin{itemize}
\item Textbook starts with CB choosing $M$.
\item Modern CB sets rates, but $M$-supply builds intuition.
\item Next: equilibrium.
\end{itemize}

\item \textbf{Money market equilibrium: where $M^s = M^d$}
\begin{itemize}
\item Given $Y$, interest rate adjusts so money demand equals supply.
\item Excess money demand $\Rightarrow$ rates rise to reduce money demand.
\item Next: turn this into LM.
\end{itemize}

\item \textbf{From money market equilibrium to the LM curve}
\begin{itemize}
\item $Y\uparrow \Rightarrow$ money demand rises; to clear, $i$ must rise.
\item Connecting clearing points gives upward-sloping LM.
\item Next: shifts.
\end{itemize}

\item \textbf{How LM shifts I: higher $P$ shifts LM left}
\begin{itemize}
\item $P\uparrow \Rightarrow M/P\downarrow$ (tighter money market).
\item To clear, $i$ higher at each $Y$ $\Rightarrow$ LM up/left.
\item Next: $M\uparrow$.
\end{itemize}

\item \textbf{How LM shifts II: higher $M^s$ shifts LM right}
\begin{itemize}
\item $M\uparrow$ increases real balances.
\item Clearing requires lower $i$ at each $Y$ $\Rightarrow$ LM down/right.
\item Next: interpretation.
\end{itemize}

\item \textbf{What LM is (and is not): avoid a common confusion}
\begin{itemize}
\item LM is a \textbf{market-clearing locus}, not a full model.
\item Need goods market too.
\item Next: IS.
\end{itemize}

\item \textbf{IS: goods market equilibrium (bridge from RBC to textbook)}
\begin{itemize}
\item IS is goods market clearing: planned spending = output (or $S=I$).
\item Lower $i$ encourages demand (esp. investment) $\Rightarrow$ $Y$ rises in equilibrium.
\item Next: visualize slope.
\end{itemize}

\item \textbf{IS curve (visual): why it slopes downward}
\begin{itemize}
\item $i\downarrow \Rightarrow$ spending/investment $\uparrow \Rightarrow$ equilibrium $Y\uparrow$.
\item IS is downward sloping in $(Y,i)$.
\item Next: shifts.
\end{itemize}

\item \textbf{Shifts in IS: what moves goods demand}
\begin{itemize}
\item Fiscal expansion / optimism / easier credit shift IS right.
\item Uncertainty / tighter credit shift IS left.
\item Next: intersection.
\end{itemize}

\item \textbf{IS--LM equilibrium: one point that clears two markets}
\begin{itemize}
\item Intersection clears goods market (IS) and money market (LM).
\item Gives short-run equilibrium $(Y,i)$ for given $M$ and $P$.
\item Next: derive AD.
\end{itemize}

\item \textbf{From IS--LM to Aggregate Demand: why AD slopes down}
\begin{itemize}
\item $P\downarrow \Rightarrow M/P\uparrow \Rightarrow$ LM right $\Rightarrow i\downarrow \Rightarrow Y\uparrow$.
\item Hence AD slopes downward in $(P,Y)$.
\item Next: policy shifts.
\end{itemize}

\item \textbf{Policy experiment I: fiscal expansion shifts AD right}
\begin{itemize}
\item $G\uparrow$ shifts IS right $\Rightarrow$ higher $Y$ at any $P$ $\Rightarrow$ AD right.
\item Output vs inflation split depends on AS (next lecture).
\end{itemize}

\item \textbf{Policy experiment II: monetary expansion shifts AD right}
\begin{itemize}
\item $M\uparrow$ shifts LM right $\Rightarrow$ higher $Y$ at any $P$ $\Rightarrow$ AD right.
\item Modern CBs do this via rates; we upgrade later.
\end{itemize}

\item \textbf{Why IS--LM is useful (and why we won’t stop here)}
\begin{itemize}
\item Useful for intuition and comparative statics.
\item Limited on expectations/credibility and modern institutions; price setting not microfounded.
\item Next: AD--AS.
\end{itemize}

\item \textbf{Preview: the AD--AS view}
\begin{itemize}
\item IS--LM gives AD; inflation requires AS.
\item AD--AS tells output vs inflation responses to shocks.
\end{itemize}

\item \textbf{Next time}
\begin{itemize}
\item Build AS and run demand vs supply shocks.
\item Start explaining inflation, not just output.
\end{itemize}
\end{enumerate}



\end{document}
